\section{Circuit Complexity --- Uniform}
%%% generated by the blueprint pipeline from TCSlib.Complexity.CircuitComplexity.Uniform
\subsection*{Circuit Complexity --- Uniform: Introduction}

Declarations of the Lean module \texttt{TCSlib.Complexity.CircuitComplexity.Uniform}.

\subsection*{Circuit Complexity --- Uniform: Results}

\begin{definition}[Prefix-free encoding condition]
\lean{BoolCircuit.PrefixFree}
\label{BoolCircuit.PrefixFree}
\leanok
Let $\alpha$ be a type and let $f$ map each element of $\alpha$ to a finite Boolean
word. The encoding satisfies the prefix-free condition if, for all $a,b\in\alpha$ and
Boolean words $s,t$, equality $f(a)s=f(b)t$ under concatenation implies both $a=b$ and
$s=t$.
\end{definition}

\begin{definition}[Recursive list encoding]
\lean{BoolCircuit.encodeList}
\label{BoolCircuit.encodeList}
\leanok
Let $\alpha$ be a type and let $f$ encode each element of $\alpha$ as a Boolean word.
Encode the empty list as the one-bit word $[\mathrm{false}]$; encode a nonempty list
with head $a$ and tail $A$ as $[\mathrm{true}]$ followed by $f(a)$ and then the encoding
of $A$.
\end{definition}

\begin{definition}[Gate-kind code]
\lean{BoolCircuit.GateKind.encode}
\label{BoolCircuit.GateKind.encode}
\leanok
\uses{BoolCircuit.GateKind}
Encode the gate kinds AND, OR, and NOT as the two-bit words $[0,0]$, $[0,1]$, and
$[1,0]$, respectively.
\end{definition}

\begin{definition}[Encoding a directed acyclic graph gate]
\lean{BoolCircuit.DAGGate.encode}
\label{BoolCircuit.DAGGate.encode}
\leanok
\uses{BoolCircuit.DAGGate, BoolCircuit.GateKind.encode, BoolCircuit.encodeList, BoolCircuit.encodeNat}
A gate labeled AND, OR, or NOT is encoded by concatenating the two-bit code for its
label with the list encoding of its argument indices, where natural numbers are encoded
in unary with a terminating zero.
\end{definition}

\begin{definition}[Encoding a directed acyclic circuit]
\lean{BoolCircuit.DAGCircuit.encode}
\label{BoolCircuit.DAGCircuit.encode}
\leanok
\uses{BoolCircuit.DAGCircuit, BoolCircuit.DAGGate.encode, BoolCircuit.encodeList, BoolCircuit.encodeNat}
For a circuit with $n$ input vertices, its encoding concatenates the unary encoding of
$n$, the list encoding of its gates, and the unary encoding of its output vertex. The
gates are listed in topological order, each gate reads only vertices with index less
than $n+i$ when it is the $i$th gate, and the output index is less than $n$ plus the
number of gates.
\end{definition}

\begin{lemma}[Prefix freeness of natural number codes]
\lean{BoolCircuit.encodeNat_prefixFree}
\label{BoolCircuit.encodeNat_prefixFree}
\leanok
\uses{BoolCircuit.PrefixFree, BoolCircuit.encodeNat}
\difficulty{4}
For natural numbers $a,b$ and Boolean lists $s,t$, if the code for $a$ followed by $s$
equals the code for $b$ followed by $t$, then $a=b$ and $s=t$.
\end{lemma}

\begin{lemma}[Prefix-free list encoding]
\lean{BoolCircuit.encodeList_prefixFree}
\label{BoolCircuit.encodeList_prefixFree}
\leanok
\uses{BoolCircuit.PrefixFree, BoolCircuit.encodeList}
\difficulty{4}
Let $A$ be a type, and let $f$ assign a finite Boolean string to each element of $A$.
Suppose $f$ is prefix-free in the strong sense that whenever $f(a)$ followed by a string
$s$ equals $f(b)$ followed by a string $t$, one has $a=b$ and $s=t$. Then the encoding
of finite lists formed from $f$ is also prefix-free in this sense.
\end{lemma}

\begin{lemma}[Prefix-free gate-kind encoding]
\lean{BoolCircuit.GateKind.encode_prefixFree}
\label{BoolCircuit.GateKind.encode_prefixFree}
\leanok
\uses{BoolCircuit.GateKind.encode, BoolCircuit.PrefixFree}
\difficulty{4}
For any gate kinds $a,b$ and Boolean lists $s,t$, if the encoding of $a$ followed by $s$
equals the encoding of $b$ followed by $t$, then $a=b$ and $s=t$.
\end{lemma}

\begin{lemma}[Prefix freeness of gate encoding]
\lean{BoolCircuit.DAGGate.encode_prefixFree}
\label{BoolCircuit.DAGGate.encode_prefixFree}
\leanok
\uses{BoolCircuit.DAGGate.encode, BoolCircuit.GateKind.encode_prefixFree, BoolCircuit.PrefixFree, BoolCircuit.encodeList_prefixFree, BoolCircuit.encodeNat_prefixFree}
\difficulty{3}
The encoding of Boolean circuit gates is prefix-free in the following sense: if the
encoding of a gate $a$ followed by a Boolean list $s$ equals the encoding of a gate $b$
followed by a Boolean list $t$, then $a=b$ and $s=t$.
\end{lemma}

\begin{theorem}[Uniqueness from circuit encoding]
\lean{BoolCircuit.DAGCircuit.eq_of_encode_eq}
\label{BoolCircuit.DAGCircuit.eq_of_encode_eq}
\leanok
\uses{BoolCircuit.DAGCircuit, BoolCircuit.DAGCircuit.encode, BoolCircuit.DAGGate.encode_prefixFree, BoolCircuit.encodeList_prefixFree, BoolCircuit.encodeNat_prefixFree}
\difficulty{3}
Let $C_1$ and $C_2$ be directed acyclic Boolean circuits, each with gates in topological
order and a valid output vertex, possibly with different numbers of inputs. If their
encodings are equal, then they have the same number of inputs, the same ordered list of
gates, and the same output vertex.
\end{theorem}

\begin{theorem}[Injectivity of circuit encoding]
\lean{BoolCircuit.DAGCircuit.encode_injective}
\label{BoolCircuit.DAGCircuit.encode_injective}
\leanok
\uses{BoolCircuit.DAGCircuit, BoolCircuit.DAGCircuit.encode, BoolCircuit.DAGCircuit.eq_of_encode_eq}
\difficulty{4}
For each natural number $n$, the Boolean-string encoding is injective on the directed
acyclic Boolean circuits with $n$ input vertices: any two such circuits with the same
encoding are equal.
\end{theorem}

\begin{lemma}[Lower bound for list encoding length]
\lean{BoolCircuit.length_encodeList_ge}
\label{BoolCircuit.length_encodeList_ge}
\leanok
\uses{BoolCircuit.encodeList}
\difficulty{4}
Let $\alpha$ be a type, $f$ a function from $\alpha$ to Boolean lists, and $as$ a finite
list of elements of $\alpha$. The encoding of $as$ has length at least $|as|+1$.
\end{lemma}

\begin{theorem}[Size bounded by encoding length]
\lean{BoolCircuit.DAGCircuit.size_le_length_encode}
\label{BoolCircuit.DAGCircuit.size_le_length_encode}
\leanok
\uses{BoolCircuit.DAGCircuit, BoolCircuit.DAGCircuit.encode, BoolCircuit.DAGCircuit.size, BoolCircuit.DAGGate.encode, BoolCircuit.encodeNat, BoolCircuit.length_encodeList_ge}
\difficulty{3}
Let $C$ be a finite Boolean circuit with $n$ input vertices, gates in topological order,
and a designated output vertex. The size of $C$ is at most the length of its encoding.
\end{theorem}

\begin{lemma}[Upper bound for list encoding length]
\lean{BoolCircuit.length_encodeList_le}
\label{BoolCircuit.length_encodeList_le}
\leanok
\uses{BoolCircuit.encodeList}
\difficulty{4}
Let $\alpha$ be a type, $f$ a function from $\alpha$ to Boolean lists, $B\in\mathbb{N}$,
and $as$ a finite list such that every $a\in as$ satisfies $|f(a)|\le B$. Then the
encoding of $as$ has length at most $|as|(B+1)+1$.
\end{lemma}

\begin{theorem}[Encoding length for fan-in-two circuits]
\lean{BoolCircuit.DAGCircuit.length_encode_le_of_isFaninTwo}
\label{BoolCircuit.DAGCircuit.length_encode_le_of_isFaninTwo}
\leanok
\uses{BoolCircuit.DAGCircuit, BoolCircuit.DAGCircuit.IsFaninTwo, BoolCircuit.DAGCircuit.encode, BoolCircuit.DAGCircuit.size, BoolCircuit.DAGGate.encode, BoolCircuit.GateKind.encode, BoolCircuit.encodeNat, BoolCircuit.length_encodeList_le}
\difficulty{5}
Let $C$ be a directed acyclic circuit with $n$ input vertices, gates in topological
order, and a valid output vertex. Suppose each gate has distinct input vertices, every
NOT gate has exactly one input, and each gate has at most two inputs. If the circuit
size is $n$ plus its number of gates, then the length of its encoding is at most $12$
times the square of its size.
\end{theorem}

\begin{definition}[Polynomial-time uniformity of a circuit family]
\lean{BoolCircuit.DAGCircuitFamily.IsPUniform}
\label{BoolCircuit.DAGCircuitFamily.IsPUniform}
\leanok
\uses{BoolCircuit.DAGCircuit.encode, BoolCircuit.DAGCircuitFamily, Complexity.PolyTimeComputable}
A family assigning to each natural number $n$ a directed acyclic circuit on $n$ input
vertices is polynomial-time uniform if there is a polynomial-time computable function on
finite Boolean lists that, on the list of $n$ true bits, returns the encoding of the
circuit assigned to $n$.
\end{definition}

\begin{theorem}[Polynomial-time uniformity implies polynomial size]
\lean{BoolCircuit.DAGCircuitFamily.IsPUniform.isPolySize}
\label{BoolCircuit.DAGCircuitFamily.IsPUniform.isPolySize}
\leanok
\uses{BoolCircuit.DAGCircuit.size_le_length_encode, BoolCircuit.DAGCircuitFamily, BoolCircuit.DAGCircuitFamily.IsPUniform, BoolCircuit.DAGCircuitFamily.IsPolySize, Complexity.PolyTimeComputable.output_length_le}
\difficulty{4}
Let $C$ be a family of directed acyclic Boolean circuits. If $C$ is
polynomial-time-uniform, then there are natural numbers $a$ and $k$ such that
$\mathrm{size}(C_n)\leq a(n+1)^k$ for every $n$.
\end{theorem}
